% Generated draft: requires mathematical and source review.
\documentclass[11pt]{article}
\usepackage[margin=25mm]{geometry}
\usepackage{fontspec}
\setmainfont{Latin Modern Roman}
\usepackage{amsmath,amssymb,mathtools,unicode-math}
\setmathfont{Latin Modern Math}
\usepackage{enumitem,needspace,xurl,hyperref}
\setlength{\parindent}{0pt}
\setlength{\parskip}{5pt}
\setlength{\emergencystretch}{3em}
\newcommand{\N}{\mathbb{N}}
\newcommand{\Z}{\mathbb{Z}}
\newcommand{\Q}{\mathbb{Q}}
\newcommand{\R}{\mathbb{R}}
\newcommand{\C}{\mathbb{C}}
\newcommand{\F}{\mathbb{F}}
\newcommand{\A}{\mathbb{A}}
\newcommand{\PP}{\mathbb{P}}
\newcommand{\E}{\mathbb{E}}
\newcommand{\eps}{\varepsilon}
\newcommand{\dd}{\,\mathrm{d}}
\begin{document}
\section*{UnsolvedMath ID: 30004789}
\textbf{Source problem code:} catalogue entry
\subsection{Definitions and mathematical statement}
The supplied record describes a conjecture about integral models of Shimura varieties with parahoric level structure. A parahoric group scheme \(\mathcal G\) alone does not specify the associated Shimura varieties, their reflex field, or the integral base. Consequently, the following is a schematic statement of the recorded requirements; the missing data and qualifications must be supplied before it becomes a fully specified conjecture.

Write \(B\) for the intended ring of reflex integers, \(\eta\) for its generic point, \(X_K\) for the intended Shimura variety at level \(K\), and \(\mathcal P_{\eta,K}\) for its asserted generic-fiber \(\mathcal G\)-shtuka. These symbols denote data presupposed by the source, not data determined by the supplied group scheme. An integral model here means a model \(\mathcal S_K\) over \(B\), together with an identification
\[
\mathcal S_K\times_B\eta\simeq X_K.
\]
Normality means that its local rings are integrally closed domains; flatness is required for its structural morphism to \(B\).

The recorded conjectural construction asks for a compatible inverse system \((\mathcal S_K)_K\) of normal, flat integral models satisfying the following conditions:
\begin{enumerate}
\item The transition morphisms for the intended changes of level are finite and \('{e}\)tale and extend the corresponding morphisms between the generic fibers.
\item The models satisfy the valuative extension property intended in the source. The supplied record does not specify the test valuation rings, the maps to be extended, the relevant level of the inverse system, or whether uniqueness is part of this property.
\item For each intended level \(K\), the generic-fiber shtuka \(\mathcal P_{\eta,K}\) extends to a \(\mathcal G\)-shtuka \(\mathcal P_K\) over \(\mathcal S_K\), compatibly with the system.
\item Every formal completion specified by the intended local uniformization assertion is identified with the corresponding integral local Shimura variety, and the extended shtuka is identified with the tautological shtuka under that identification. The record does not specify the completion loci, local Shimura data, base changes, framing choices, or precise compatibility morphisms.
\end{enumerate}

The permitted levels, the integral base, the definition of a \(\mathcal G\)-shtuka in use, and the precise valuative and formal-local assertions remain unresolved. In particular, the phrase ``every formal completion'' cannot be given a precise quantifier from this record alone. The requirements above preserve the supplied target without asserting a more specific extension theorem or uniformization theorem.
\subsection{Short English statement}
Construct compatible normal, flat integral models of the intended Shimura varieties with finite \('{e}\)tale transition maps, an extension property, extended universal shtukas, and compatible descriptions of their formal completions by integral local Shimura varieties. The supplied record omits data needed to state these requirements precisely. Its partially solved label and verification claims are supplied curation, not independent mathematical verification.
\subsection{Sources}
\begin{itemize}
\item Local UnsolvedMath problem registry: permanent ID \(30004789\), source problem code \texttt{catalogue entry}, title ``Integral Shimura Models with Universal Shtukas.'' The supplied record gives proposed year 2021 and names no proposer.
\item Supplied report citation: \emph{Automorphic Forms, Geometry and Arithmetic (2022)}, Oberwolfach Reports, pp.~2089--2156, DOI \texttt{10.4171/owr/2021/39}. \url{https://doi.org/10.4171/owr/2021/39}. This is also the supplied literature-source link labelled ``original Oberwolfach report.''
\item Supplied literature source: Pappas--Rapoport, ``integral local Shimura varieties.'' \url{https://doi.org/10.1017/S1474748025101345}.
\item Supplied literature source: Pappas--Rapoport, preprint. \url{https://arxiv.org/abs/2204.02829}.
\item Literature assessment dated exactly \texttt{2026-08-22}: the supplied curation labels the problem partially solved and reports representability and local-model assertions for abelian-type data under stated prime and group hypotheses, extending PEL cases, while describing arbitrary parahoric data as not fully settled. Those prime and group hypotheses are not reproduced in the record. This is supplied dated curation, not independently verified; its time-sensitive status and cited claims require independent checking.
\item The registry labels the statement \texttt{corrected\_verified} and says its standalone statement and hypotheses were checked against the official report. These are supplied source-verification claims; no independent source consultation is asserted here.
\end{itemize}
\end{document}
